\documentclass[11pt,letterpaper]{article}
\usepackage[margin=1in]{geometry}
\usepackage{amsmath,amssymb,amsthm,mathtools}
\usepackage{enumitem}
\usepackage[hidelinks]{hyperref}
\emergencystretch=3em

\newtheorem{theorem}{Theorem}[section]
\newtheorem{lemma}[theorem]{Lemma}
\newtheorem{corollary}[theorem]{Corollary}
\newtheorem{proposition}[theorem]{Proposition}
\newtheorem{claim}{Claim}
\theoremstyle{definition}
\newtheorem{definition}[theorem]{Definition}
\newtheorem{remark}[theorem]{Remark}

\newcommand{\LP}{\mathrm{LP}}
\newcommand{\OPT}{\mathrm{OPT}}
\newcommand{\R}{\mathbb{R}}
\newcommand{\Z}{\mathbb{Z}}
\newcommand{\load}{\mathrm{load}}
\newcommand{\eps}{\varepsilon}

\title{A Near-Logarithmic Approximation for\\ All-or-Nothing Multicommodity Flow}
\author{Zijian Zeng\\ UCSI University\\ \texttt{1002266693@ucsiuniversity.edu.my}}
\date{}

\begin{document}
\maketitle

\begin{abstract}
In the all-or-nothing multicommodity flow problem we are given an undirected graph with unit
edge capacities and $k$ terminal pairs, and we seek a largest set of pairs that can
simultaneously send one unit of (splittable) flow each. The best known approximation ratio,
due to Chekuri, Khanna and Shepherd (STOC 2005), is $O(\log^2 k)$, and Bansal, Jambulapati and
Saranurak (ICALP 2026) conjectured that a $\log^{1+o(1)} k$ approximation exists. We confirm
this conjecture: we give a polynomial-time algorithm that routes at least
$\LP/(\log k\cdot 2^{O(\sqrt{\log\log k})})$ pairs with congestion one, where $\LP$ is the value
of the natural multicommodity flow relaxation. Hence the integrality gap of this relaxation is
$\log^{1+o(1)}k$; the known lower bound is $\log^{1/2-\delta}n$.

The proof has two parts. First, a single decomposition computed by the integral-node-weight
flow-expander decomposition algorithm of Bansal, Jambulapati and Saranurak, with node weights obtained
by discretizing the fractional solution at granularity $\LP/(2k)$, yields the well-linked
decomposition of Chekuri, Khanna and Shepherd with loss $\log^{1+o(1)}k$ instead of
$O(\log^2k)$. This is the decomposition asked for by Bansal, Jambulapati and Saranurak, and
the argument is short. Second, we give a complete and self-contained proof of the
congestion-one rounding inside flow-linked instances that the decomposition requires. This
rounding was stated by Chekuri, Khanna and Shepherd in 2005 together with a remark that its
clustering had to be strengthened; the strengthened argument appeared in their 2013 journal
paper, phrased for R\"acke-tree instances. Our proof follows their architecture, replaces
cluster centers by a min-cut criterion, and uses a degree reduction to make clusters
vertex-disjoint. We also extend the rounding to non-uniform paired weights.
\end{abstract}

\section{Introduction}

\paragraph{The problem.}
Let $G=(V,E)$ be an undirected graph with unit edge capacities and let
$(s_1,t_1),\dots,(s_k,t_k)$ be (indexed) terminal pairs. A set $I\subseteq[k]$ is
\emph{routable} if there is a multicommodity flow that sends one unit from $s_i$ to $t_i$ for
every $i\in I$ and whose total load on every edge is at most one; flow of a commodity may be
split along many paths. The \emph{all-or-nothing multicommodity flow problem} (AN-MCF) asks
for a routable set of maximum cardinality. It was introduced by Chekuri, Mydlarz and
Shepherd~\cite{CMS07} and studied in general graphs by Chekuri, Khanna and
Shepherd~\cite{CKS04,CKS13}, as a relaxation of the edge-disjoint paths problem (EDP) that,
unlike EDP, admits polylogarithmic approximation ratios: the natural multicommodity flow
relaxation has an $\Omega(\sqrt n)$ integrality gap for EDP but only a polylogarithmic one for
AN-MCF. Call $x\in[0,1]^k$ \emph{fractionally routable} if some multicommodity flow with load
at most one on every edge routes $x_i$ units between $s_i$ and $t_i$ for every $i$. The
relaxation is
\begin{equation}\label{eq:lp}
\LP=\max\Bigl\{\textstyle\sum_{i=1}^k x_i\;:\;x\in[0,1]^k\text{ fractionally routable}\Bigr\};
\end{equation}
it has a polynomial-size edge formulation, and clearly $\LP\ge\OPT$.

\paragraph{Prior work.}
Chekuri, Khanna and Shepherd~\cite{CKS04,CKS13} gave the first polylogarithmic
approximation, $O(\log^3 n\log\log n)$, using R\"acke's hierarchical
decompositions~\cite{Racke02,HHR03}. In~\cite{CKS05} they made precise the notion of a
\emph{well-linked decomposition} (used earlier in~\cite{CKS04} and in work on planar EDP) and
introduced its weighted version: a fractional solution is converted into a collection of
vertex-disjoint instances in each of which the terminals are fractionally flow-linked, losing
a factor $O(\beta\log k)$, where $\beta=O(\log k)$ is the flow-cut gap for product
multicommodity flows~\cite{LR99,LLR95,AR98}; inside a flow-linked instance, a congestion-one
rounding loses only a constant~\cite[Theorem~3.1]{CKS05}. This gives the current best ratio
$O(\log^2 k)$ for general graphs~\cite[Theorem~3.2]{CKS05} (and $O(\log k)$ in planar graphs).
On the negative side, AN-MCF is hard to approximate within $\log^{1/2-\delta}n$ for every
$\delta>0$ unless $\mathrm{NP}\subseteq\mathrm{ZPTIME}(n^{\mathrm{polylog}\, n})$, and the
relaxation~\eqref{eq:lp} has integrality gap $\log^{1/2-\delta}n$~\cite{ACGKTZ10}; the problem is
APX-hard already on trees~\cite{GVY97}.

Recently Bansal, Jambulapati and Saranurak~\cite{BJS26} gave the first polynomial-time
flow-expander decompositions with overhead $\log^{1+o(1)}$, circumventing the flow-cut gap. In
their list of open questions~\cite[Section~5]{BJS26} they observe that the $O(\log^2k)$ bound
for AN-MCF is the product of the flow-cut gap and the recursion depth of the well-linked
decomposition, and that their technique suggests these might be telescoped; but they point
out that the connection is not black-box, since the well-linked decomposition of~\cite{CKS05}
interleaves concurrent-flow computations with maximum-throughput flow steps. They ask whether
the decomposition of~\cite{CKS05} can be built with $\log^{1+o(1)}k$ overhead, note that this
would immediately imply the following conjecture, and state:

\begin{quote}
\textbf{Conjecture 5.1 of~\cite{BJS26}.} There is a polynomial-time $\log^{1+o(1)}k$-approximation
algorithm for the all-or-nothing flow problem.
\end{quote}

\subsection{Our results}

All logarithms are binary. We fix, once and for all, a nondecreasing function
$\Gamma:\Z_{\ge 0}\to\R_{\ge1}$ with
\begin{equation}\label{eq:gamma}
\Gamma(a)=O\bigl(\log(a+2)\cdot\exp\bigl(O(\sqrt{\log\log(a+4)})\bigr)\bigr)
\end{equation}
such that the decomposition theorem of~\cite{BJS26} holds with cut bound $\phi|A|\Gamma(|A|)$
(Theorem~\ref{thm:bjs} and the discussion after it).

\begin{theorem}[Main theorem]\label{thm:main}
There is a polynomial-time algorithm that, given an undirected graph with unit edge
capacities and $k$ terminal pairs, outputs a routable set of pairs together with a routing of
congestion one, of size at least
\[
\frac{\LP}{c\,\Gamma(4k)}\;\ge\;\frac{\OPT}{C\log(k+2)\exp\bigl(C\sqrt{\log\log (k+4)}\bigr)},
\]
where $c,C$ are absolute constants. In particular the integrality gap of~\eqref{eq:lp} is
$\log^{1+o(1)}k$.
\end{theorem}

This confirms Conjecture~5.1 of~\cite{BJS26}. The guarantee depends only on $k$ (not on the
number of vertices), no congestion is violated, and shared endpoints, coincident endpoints and
repeated pairs are allowed. Apart from the call to the decomposition algorithm of~\cite{BJS26},
the algorithm is deterministic. Theorem~\ref{thm:main} is the composition of the following two
results.

\paragraph{Well-linked decomposition with $\log^{1+o(1)}k$ loss.}
Following~\cite{CKS05}, for a set $X$ of vertices and weights $\pi:X\to\R_{\ge0}$, $X$ is
\emph{$\pi$-flow-linked} in $G$ if the product demand $\pi(u)\pi(v)/\pi(X)$ between every two
distinct $u,v\in X$ is routable with congestion one.

\begin{theorem}[see Theorem~\ref{thm:decomp}]\label{thm:decomp-intro}
Given a feasible fractional multicommodity flow of total value $F$ for $k$ pairs with distinct
degree-one endpoints, one can compute in polynomial time vertex-disjoint induced subgraphs
$G_1,\dots,G_\ell$ and, for each $j$, a set $P_j$ of pairs inside $G_j$ with weights
$\pi_j$ on their endpoints such that $\pi_j(s_i)=\pi_j(t_i)\le1$, the endpoints $X_j$ of
$P_j$ are $\pi_j$-flow-linked in $G_j$, and $\sum_j\pi_j(X_j)\ge F/(c\,\Gamma(4k))$.
\end{theorem}

This is the well-linked decomposition of~\cite[Theorem~2.1]{CKS05} with the loss
$O(\log^2 k)$ replaced by $\log^{1+o(1)}k$, i.e.\ the decomposition asked for
in~\cite[Section~5]{BJS26}. As anticipated in~\cite{BJS26}, Theorem~\ref{thm:main} follows by
combining it with congestion-one rounding inside flow-linked instances; we now discuss this
second step, which is where most of the work of this paper lies.

\paragraph{Congestion-one rounding, with a complete proof.}
The rounding step is the following statement~\cite[Theorem~3.1]{CKS05}: if $X$ is
$\alpha$-flow-linked (i.e.\ $\pi\equiv\alpha$), then for every matching $M$ on $X$ one can
route $\Omega(\alpha|M|)$ pairs of $M$ with congestion one. In~\cite{CKS05} this theorem is
attributed to the conference version~\cite{CKS04} of~\cite{CKS13}, with a footnote stating
that the clustering of~\cite{CKS04} needs to be strengthened for congestion one and that
details would appear in the journal version. The journal version~\cite[Section~3.4]{CKS13}
contains the strengthened clustering. Its theorem~\cite[Theorem~3.8]{CKS13} is stated for
``nice'' instances defined through a R\"acke tree, which is used to obtain a flow to a common
distribution; the general case of clusters sharing vertices is indicated briefly at the end
of the section, and the weighted case is outlined. Since the rounding is the step where
congestion exactly one is obtained, we give a complete self-contained proof for arbitrary
flow-linked sets (Section~\ref{sec:rounding}).

\begin{theorem}[Theorem~\ref{thm:round} and Corollary~\ref{cor:paired}]\label{thm:round-intro}
\begin{enumerate}[label=(\alph*),nosep]
\item If $X$ is $\alpha$-flow-linked in $G$, $0<\alpha\le1$, and $M$ is a matching on $X$,
then one can compute in polynomial time a set of at least $\max\{1,\alpha|M|/2^{15}\}$ pairs of
$M$ that is routable with congestion one.
\item If the vertices of $X$ have degree one, $M$ is a perfect matching on $X$,
$\pi:X\to[0,1]$ satisfies $\pi(s)=\pi(t)$ for $st\in M$, and $X$ is $\pi$-flow-linked, then
one can compute in polynomial time a set of at least $\max\{1,\pi(X)/2^{19}\}$ pairs of $M$
routable with congestion one.
\end{enumerate}
\end{theorem}

Part (b), a ``non-uniform'' rounding, follows from part (a) by the copying device
of~\cite[Lemma~1.1]{CKS05}, in which the unit edges at the original terminals act as admission
bottlenecks.

\paragraph{Other consequences and limits.}
The reduction behind Theorem~\ref{thm:decomp-intro} accepts any flow-expander decomposition
algorithm for integral node weightings. Anand~\cite{Anand26} recently improved flow-expander
decomposition to overhead $O(\log n\log\log n)$. His bound in terms of the total weight $|A|$
requires every vertex to have weight at least one, which forces $|A|\ge n$, so we were not able
to obtain an $O(\log k\log\log k)$ ratio. If the node-weighted bound stated in Appendix~C
of~\cite{Anand26} (zero weights allowed, overhead $O(\log n\log\log n)$) is used, the same
reduction gives an $O(\log n\log\log n)$-approximation (Remark~\ref{rem:anand}); this is
conditional on that statement and is not a bound in terms of $k$. For the weighted (profit)
version we only obtain $\log^{2+o(1)}k$ (Remark~\ref{rem:weighted}); a $\log^{1+o(1)}k$ bound
for it, and the node-capacitated version, remain open (Section~\ref{sec:discussion}).

\subsection{Technical overview}

\paragraph{The decomposition.}
The decomposition of~\cite{CKS05} fixes a path decomposition of the fractional solution and
recursively cuts sparse cuts with respect to the current weights $\gamma(u,H)$, the flow of
paths fully contained in the current piece $H$. These weights shrink whenever a cut destroys
paths, so the weighting changes during the recursion, whereas the decomposition
of~\cite{BJS26} works with one fixed integral node weighting $A$ and charges the cut to $|A|$.
We therefore fix the weights in advance. We discretize the fractional solution $x$ at
granularity $\eps=F/(2k)$: both endpoints of pair $i$ receive the integral weight
$d_i=\lfloor x_i/\eps\rfloor$ (zero if $x_i<\eps$), in the spirit of the uniformization
of~\cite[Lemma~1.1]{CKS05}. Then $K=\sum_id_i\le2k$, so $|A|\le4k$, while $\eps K\ge F/4$. We
apply~\cite{BJS26} once (up to a constant number of trials to calibrate an unknown constant),
with $\phi=\Theta(\eps/\Gamma(4k))$, so that at most $\eps K/4$ edges are cut. Two simple
observations finish the proof. First, a pair whose endpoints are separated sends all of its
flow $x_i\ge\eps d_i$ across the cut, so separated pairs carry at most a quarter of the total
weight; this is the charging of~\cite[Lemma~2.2]{CKS05}, now against a fixed weighting. Second,
flow expansion with respect to $A$ implies flow-linkedness of every sub-collection of weighted
terminals, so after discarding separated pairs the remaining endpoints in each component are
$\phi A$-flow-linked, and no weight is ever recomputed.

\paragraph{Admission.}
The pieces carry non-uniform weights $\phi d_i$. To use the uniform rounding we re-discretize
and attach pendant copies of each terminal; since the original terminals have degree one, two
copies of the same pair would have to share the unit edge at $s_i$, so a congestion-one routing
of copy pairs uses at most one copy of each pair, and the rounding does not need to know about
copies. This gives Theorem~\ref{thm:round-intro}(b).

\paragraph{The rounding.}
Our proof of Theorem~\ref{thm:round-intro}(a) follows the architecture of~\cite[Section~3.4]{CKS13}:
terminals are grouped into clusters of $\Theta(1/\alpha)$ terminals, a sparse set of pairs is
selected, every selected terminal spreads one unit to the terminals of its (possibly
enlarged) cluster using at most half of the capacity of every edge except bridges, the
received flow is sent to a common distribution using the other half of the capacity, and the
excess on bridges is removed by an integrality argument~\cite[Lemma~3.10]{CKS13}. We make
three changes that allow a short and complete proof: (i) a max-flow/min-cut argument
(Lemma~\ref{lem:criterion}) gives a simple sufficient condition for a terminal to spread its
unit, which removes the need for cluster centers; (ii) a flow-preserving degree reduction makes
all clusters vertex-disjoint, which removes the case analysis caused by shared connector
vertices; (iii) the common-distribution flow is obtained directly from flow-linkedness rather
than from a R\"acke tree.

\paragraph{On the contribution.}
Once Theorem~3.1 of~\cite{BJS26} is available, the decomposition (Theorem~\ref{thm:decomp-intro})
is a short argument, and~\cite{BJS26} anticipated that such a decomposition would imply their
conjecture. The longer part of the paper is the congestion-one rounding; its structure follows
\cite[Section~3.4]{CKS13}, and our contribution there is a complete proof for arbitrary
flow-linked sets, with the simplifications above, and the non-uniform extension. We view the
paper as settling the conjecture of~\cite{BJS26} by a short reduction together with a careful,
self-contained treatment of the rounding step that the reduction needs.

\paragraph{Organization.}
Section~\ref{sec:prelim} gives definitions and the tools we use. Section~\ref{sec:rounding}
proves the rounding theorem, Section~\ref{sec:decomp} the decomposition theorem, and
Section~\ref{sec:main} the main theorem and its corollaries. Section~\ref{sec:discussion}
discusses lower bounds and open problems.

\section{Preliminaries}\label{sec:prelim}

\paragraph{Graphs and flows.}
Graphs are undirected and may have parallel edges; every edge has capacity one. For a graph
$H$ and $Z\subseteq V(H)$, $\delta_H(Z)$ is the multiset of edges with exactly one endpoint in
$Z$, and $H[Z]$ is the induced subgraph. An edge $e$ of $H$ is a \emph{bridge of $H$} if
$H-e$ has more connected components than $H$; an edge with a parallel copy is never a bridge.
If $e$ is a bridge of a connected graph $H$, then $H-e$ has exactly two components.
A flow from $u$ to $v$ of value $\nu$ is a nonnegative function on the arcs of the
bidirected version of $H$ satisfying conservation at all vertices other than $u,v$, with net
outflow $\nu$ at $u$. The load of a collection of flows on an edge $e$ is the sum of all
their values on the two arcs of $e$. A demand $D$ (a nonnegative function on unordered pairs
of distinct vertices) is \emph{routable with congestion $\kappa$} if there are flows of value
$D(u,v)$ between $u$ and $v$ for all pairs with total load at most $\kappa$ on every edge. A
set of terminal pairs is \emph{routable} if the demand placing one unit on each pair is
routable with congestion one. Routability of a given set is decidable by linear programming.
If $Z$ contains $u$ but not $v$, then every flow from $u$ to $v$ of value $\nu$ has load at
least $\nu$ on $\delta_H(Z)$.

\begin{definition}[flow-linked sets~\cite{CKS05}]
Let $X\subseteq V(H)$ and $\pi:X\to\R_{\ge0}$ with $\pi(X)>0$. Then $X$ is
\emph{$\pi$-flow-linked in $H$} if the demand $\pi(u)\pi(v)/\pi(X)$ on every pair of distinct
$u,v\in X$ is routable in $H$ with congestion one. For a constant $\alpha$, $X$ is
\emph{$\alpha$-flow-linked} if it is $\pi$-flow-linked for $\pi\equiv\alpha$, i.e.\ the
demand $\alpha/|X|$ on every pair is routable with congestion one.
\end{definition}

If $X$ is $\pi$-flow-linked in $H$ and $|X\cap\mathrm{supp}(\pi)|\ge2$, then all vertices of
$X$ with positive weight lie in one component of $H$.

\begin{definition}[flow expansion~\cite{BJS26}]
Let $A:V(H)\to\R_{\ge0}$. A demand $D$ \emph{respects} $A$ if
$\sum_{y\ne x}D(x,y)\le A(x)$ for every vertex $x$. The weighting $A$ is
\emph{$\phi$-flow-expanding in $H$} if every demand that respects $A$ is routable in $H$ with
congestion $1/\phi$. For $U\subseteq V(H)$,
$A|_U$ denotes the restriction of $A$ to $U$.
\end{definition}

\begin{theorem}[{Bansal, Jambulapati and Saranurak~\cite[Theorem~3.1]{BJS26}}]\label{thm:bjs}
For every undirected graph $H$, every integral node weighting $A:V(H)\to\Z_{\ge0}$ and every
$\phi>0$, one can compute in polynomial time a set $C\subseteq E(H)$ such that
$|C|\le\phi|A|\log(|A|)\exp(O(\sqrt{\log\log|A|}))$ and, for each connected component $U$ of
$H-C$, the weighting $A|_U$ is $\phi$-flow-expanding in $H[U]$.
\end{theorem}

We use this theorem as a black box, with the following conventions.
\begin{itemize}[nosep]
\item \emph{The overhead function.} We write the bound as $|C|\le\phi|A|\Gamma(|A|)$ with
$\Gamma=c_0\Gamma_1$, where $\Gamma_1(a)=\lceil\log(a+4)\rceil\cdot2^{\lceil\sqrt{\log\log(a+16)}\,\rceil}$
is explicit and $c_0\ge1$ is the absolute constant implicit in~\cite{BJS26}. Our algorithm
does not need to know $c_0$ (see the proof of Theorem~\ref{thm:decomp}). For $|A|\le1$ every
respecting demand is zero (a respecting demand on $\{x,y\}$ is at most $\min\{A(x),A(y)\}$), so
$C=\emptyset$ suffices. The internal parameters of the algorithm in~\cite{BJS26} are meant for
large $|A|$; for $|A|$ bounded by an absolute constant one may instead use classical
cut-and-recurse, which with the $O(\log|A|)$ flow-cut gap for product
demands~\cite{LR99,LLR95,AR98} has overhead $O(\log^2|A|)=O(1)$. We take $c_0$ large enough to
cover these cases.
\item \emph{A factor of two.} The stopping rule of the algorithm in~\cite{BJS26} combined with
their Fact~2.1 literally certifies $(\phi/2)$-flow expansion; if only this is available, one
calls the algorithm with $2\phi$, which changes $c_0$ by a factor two. We use the theorem as
stated.
\item \emph{Multigraphs.} Parallel edges can be subdivided without changing any flow or cut
quantity.
\item \emph{Randomness.} The algorithm of~\cite{BJS26} is built from linear programming, sweep
cuts and the sparse neighborhood cover of their Lemma~3.2; we have not verified whether the
latter is deterministic, and we use no property of the algorithm beyond the statement above.
All other steps of our algorithms are deterministic.
\end{itemize}

The following fact connects flow expansion to flow-linkedness.

\begin{lemma}\label{lem:monotone}
Let $A$ be $\phi$-flow-expanding in $H$ and let $Y\subseteq V(H)$ with $A(Y)>0$. Then $Y$ is
$\pi$-flow-linked in $H$ for $\pi=\phi A|_Y$.
\end{lemma}
\begin{proof}
The $\pi$-product demand $D(u,v)=\phi A(u)A(v)/A(Y)$ ($u\ne v\in Y$) has marginal at most
$\phi A(u)(A(Y)-A(u))/A(Y)\le\phi A(u)$ at $u\in Y$ and is zero elsewhere. So $D/\phi$ respects
$A$ and is routable with congestion $1/\phi$; scaling by $\phi$ routes $D$ with congestion
one.
\end{proof}

The next lemma is~\cite[Lemma~3.10]{CKS13}; we include its short proof.

\begin{lemma}[bridge integrality]\label{lem:bridge}
Let $H$ be a graph and $N$ a set of terminal pairs. Suppose flows of value one for all pairs
in $N$ have total load at most one on every non-bridge of $H$ and strictly less than two on
every bridge of $H$. Then $N$ is routable in $H$ with congestion one, and such a routing can
be computed in polynomial time.
\end{lemma}
\begin{proof}
For each pair, decompose its flow into flows along paths and cycles and discard the cycles;
loads do not increase, and each pair is now routed along simple paths. If $e$ is a bridge,
a simple path from $s$ to $t$ uses $e$ exactly once if $e$ separates $s$ from $t$ and not at
all otherwise. Hence the load on $e$ is the number of pairs of $N$ separated by $e$, an
integer smaller than two, i.e.\ at most one.
\end{proof}

\paragraph{Normalized instances.}
An AN-MCF instance is \emph{normalized} if its $2k$ terminals are distinct vertices of degree
one. Every instance can be normalized: for every pair $i$ add two new vertices $s_i',t_i'$ and
unit edges $s_is_i'$, $t_it_i'$, and replace the pair by $(s_i',t_i')$. Unit flows extend
along the new edges (each is used only by its own pair) and project back by deleting them,
so routable sets, $\OPT$ and $\LP$ are unchanged (a pair with $s_i=t_i$ becomes a pair at
distance two). Parallel edges of the input may be subdivided. Hence we assume throughout that
the input is normalized and simple.

\section{Congestion-one rounding in flow-linked instances}\label{sec:rounding}

\begin{theorem}\label{thm:round}
Let $G$ be a connected graph, $X\subseteq V(G)$ a set of $q\ge2$ vertices, $M$ a matching on
$X$ (a non-empty set of pairwise disjoint pairs of elements of $X$), and $0<\alpha\le1$, such
that $X$ is $\alpha$-flow-linked in $G$. Then one can compute in polynomial time a set $M'\subseteq M$
with $|M'|\ge\max\{1,\alpha|M|/2^{15}\}$ and a routing of $M'$ with congestion one.
\end{theorem}

\begin{corollary}\label{cor:paired}
Let $G$ be a graph, $M$ a set of pairs whose endpoint set $X$ consists of distinct vertices of
degree one, and $\pi:X\to[0,1]$ with $\pi(s)=\pi(t)$ for every $st\in M$ and $\pi(X)>0$, such
that $X$ is $\pi$-flow-linked in $G$. Then one can compute in polynomial time a set
$M'\subseteq M$ with $|M'|\ge\max\{1,\pi(X)/2^{19}\}$ and a routing of $M'$ with congestion
one.
\end{corollary}

As discussed in the introduction, Theorem~\ref{thm:round} is~\cite[Theorem~3.1]{CKS05} (with
an explicit constant); the proof below follows the strategy of~\cite[Section~3.4]{CKS13} but
is self-contained. Terminals not covered by $M$ are allowed; they serve only as receivers of
flow. Sections~\ref{sec:reductions}--\ref{sec:selection} prove
Theorem~\ref{thm:round}, and Section~\ref{sec:paired} proves Corollary~\ref{cor:paired}.

\subsection{Two reductions}\label{sec:reductions}

\begin{lemma}\label{lem:reduce}
Given $(G,X,M,\alpha)$ as in Theorem~\ref{thm:round}, one can construct in polynomial time a
connected graph $\hat G$ and a set $\hat X\subseteq V(\hat G)$ with a bijection
$x\mapsto\hat x$ from $X$ to $\hat X$ such that (a) every vertex of $\hat X$ has degree one
in $\hat G$; (b) every vertex of $\hat G$ has at most three distinct neighbours; (c) $\hat X$
is $\alpha$-flow-linked in $\hat G$; and (d) if a set of pairs $\{\hat s\hat t\}$ is routable
in $\hat G$, the corresponding set $\{st\}$ is routable in $G$, and the routing can be
transformed in polynomial time.
\end{lemma}
\begin{proof}
First attach to every $x\in X$ a new vertex $x'$ by a unit edge and let $X'=\{x'\}$. The
product flow between $u,v\in X$ extends along $u'u$ and $vv'$; the edge $x'x$ then carries
$\alpha(q-1)/q<1$, so $X'$ is $\alpha$-flow-linked. Routings project back by deleting the
pendant edges.

Next (after deleting loops, which carry no useful flow), let $v$ be a vertex not in $X'$
with incident edges $e_1,\dots,e_d$, $d\ge4$ (counted with multiplicity). Replace $v$ by a path $v_1,\dots,v_d$, attach $e_i$ to $v_i$, and join each
$v_l$ and $v_{l+1}$ by $d$ parallel edges. Do this for every such vertex. Every new vertex has
at most three distinct neighbours, vertices of degree at most three are unchanged, and the
graph stays connected. A flow in the new graph projects to the old one by contracting the
paths; loads on original edges are unchanged. Conversely, consider flows in the old graph
between vertices of $X'$ with load at most one on every edge, decomposed into simple paths.
A path through $v$ enters through some $e_i$ and leaves through some $e_j$; reroute it along
$v_i,\dots,v_j$ and split its value evenly among the $d$ parallel edges of each step. The
paths crossing the step between $v_l$ and $v_{l+1}$ carry at most
$\sum_i\load(e_i)\le d$ in total, so each parallel edge carries at most one. Applying this at
every split vertex lifts any routing with congestion one; in particular the product flow,
so (c) holds. Vertices of $X'$ have degree one and are not split, so (a) holds.
\end{proof}

By Lemma~\ref{lem:reduce} we assume, until the end of Section~\ref{sec:selection}, that $G$
is connected, every terminal (vertex of $X$) has degree one, every vertex has at most three
distinct neighbours, and $X$ is $\alpha$-flow-linked. ``Bridge'' always means a bridge of
this graph $G$. We set
\[
\rho:=\alpha/4 .
\]

\subsection{Common distribution and assembly}\label{sec:assembly}

\begin{lemma}[common distribution]\label{lem:common}
There are flows $h_v$ ($v\in X$), where $h_v$ sends $\rho/q$ units from $v$ to every
$w\in X\setminus\{v\}$, with total load at most $1/2$ on every edge.
\end{lemma}
\begin{proof}
Let $\psi_{vw}$ be flows routing $\alpha/q$ between every pair $\{v,w\}$ of distinct
terminals with total load at most one. Let $h_v$ consist of $\frac14\psi_{vw}$, oriented from
$v$ to $w$, for all $w\ne v$; it sends $\alpha/(4q)=\rho/q$ to $w$. Each $\psi_{vw}$ is used
with factor $\frac14$ by $h_v$ and by $h_w$, so the total load is at most half that of all
$\psi_{vw}$.
\end{proof}

\begin{definition}
For $s\in V(G)$ and $\gamma:X\to\R_{\ge0}$ with $\sum_v\gamma(v)=1$, a \emph{flow from $s$ to
$\gamma$} is a function on arcs whose net outflow at every vertex $y$ equals
$[y=s]-\gamma(y)$ (with $\gamma(y)=0$ for $y\notin X$). Let $S\subseteq X$. A \emph{spreading
system} for $S$ consists of flows $g_s$ from $s$ to $\gamma_s$, $s\in S$, such that
$\sum_{s\in S}\gamma_s(v)\le\rho$ for every $v\in X$, and the total load of all $g_s$ is at
most $1/2$ on every non-bridge and at most one on every bridge.
\end{definition}

\begin{lemma}[assembly]\label{lem:assembly}
If $M'\subseteq M$ and the set $S$ of endpoints of $M'$ has a spreading system, then $M'$ is
routable with congestion one, and a routing can be computed in polynomial time.
\end{lemma}
\begin{proof}
For $s\in S$ let $\sigma_s:=g_s+\sum_{v\in X}(\gamma_s(v)/\rho)\,h_v$. The term
$(\gamma_s(v)/\rho)h_v$ moves $\gamma_s(v)/q$ from $v$ to every $w\ne v$ and leaves
$\gamma_s(v)/q$ at $v$, so $\sigma_s$ is a flow from $s$ to the uniform distribution on $X$: the net inflow at
$w\in X$ is $\gamma_s(w)/q+\sum_{v\ne w}\gamma_s(v)/q=1/q$. For $st\in M'$, the
sum of $\sigma_s$ and the reversal of $\sigma_t$ is therefore a flow from $s$ to $t$ of value
one. The total load of all these flows is at most the load of $\sum_s g_s$ plus
$\sum_{v\in X}\bigl(\sum_{s}\gamma_s(v)/\rho\bigr)\load(h_v)\le\sum_v\load(h_v)\le 1/2$ by
Lemma~\ref{lem:common}. Thus the load is at most $1$ on non-bridges and at most $3/2<2$ on
bridges, and Lemma~\ref{lem:bridge} gives a routing with congestion one.
\end{proof}

\subsection{A criterion for spreading}\label{sec:criterion}

\begin{definition}
Let $Q$ be a connected subgraph of $G$, $s\in V(Q)$ and $R\subseteq X\cap V(Q)$. We say that $s$
\emph{spreads in $(Q,R)$} if there is a flow in $Q$ from $s$ to some $\gamma$ with
$\gamma(v)\le\rho$ for $v\in R$ and $\gamma(v)=0$ for $v\notin R$, whose load is at most $1/2$
on every edge of $Q$ that is not a bridge of $G$ and at most one on every edge of $Q$ that is a
bridge of $G$.
\end{definition}

(Up to normalization, this is the notion of a \emph{pseudocenter} in~\cite{CKS13}.)

\begin{lemma}[criterion]\label{lem:criterion}
Let $Q$ be a connected subgraph of $G$, $s\in V(Q)$ and $R\subseteq X\cap V(Q)$ with
$\rho|R|\ge1$. Suppose that for every edge $e$ of $Q$ that is a bridge of $Q$ but not a bridge
of $G$, the component of $Q-e$ containing $s$ contains at least $1/\rho$ vertices of $R$.
Then $s$ spreads in $(Q,R)$, and the flow can be computed by a maximum-flow computation.
\end{lemma}
\begin{proof}
Give every edge of $Q$ capacity one if it is a bridge of $G$ and $1/2$ otherwise, add a sink
$z$ and an arc of capacity $\rho$ from every $v\in R$ to $z$. It suffices to show that every
$s$--$z$ cut has capacity at least one: then a maximum $s$--$z$ flow, scaled to value one,
is the required flow, with $\gamma(v)$ the flow on the arc $vz$. Such a cut is given by $Z\subseteq V(Q)$ with $s\in Z$,
and its capacity is $c(\delta_Q(Z))+\rho|R\cap Z|$. If $\delta_Q(Z)=\emptyset$ then $Z=V(Q)$
and the capacity is at least $\rho|R|\ge1$. If $|\delta_Q(Z)|\ge2$, then
$c(\delta_Q(Z))\ge1$. If $\delta_Q(Z)=\{e\}$, then $e$ is a bridge of $Q$ and $Z$ is the
component of $Q-e$ containing $s$; either $e$ is a bridge of $G$ and $c(e)=1$, or by
hypothesis $\rho|R\cap Z|\ge1$.
\end{proof}

\subsection{Clusters, bad bridges and tags}\label{sec:clusters}

\begin{lemma}[clusters]\label{lem:clusters}
Let $W\ge1$ with $\rho q\ge W$. One can compute a partition $V(G)=V_1\cup\dots\cup V_p$ such
that every $G[V_j]$ is connected and $W\le\rho|X\cap V_j|<4W$ for every $j$.
\end{lemma}
\begin{proof}
Since $\rho q\ge W\ge1$ we have $q\ge 4$, so $G$ is not a single edge, and the neighbour of a
terminal is a non-terminal. Let $T$ be a spanning tree rooted at a non-terminal $r$; every
vertex has at most three children and every terminal is a leaf of $T$ (it has degree one).
Write $\rho(Y)=\rho|X\cap Y|$. Maintain a subtree $T'$ containing $r$, initially $T$. While
$\rho(V(T'))\ge W$, pick a vertex $x$ of $T'$ whose subtree $T'_x$ in $T'$ satisfies
$\rho(V(T'_x))\ge W$ while $\rho(V(T'_y))<W$ for every child $y$ of $x$ (walk down from $r$).
Since $\rho<1\le W$, $x$ is not a leaf, so it is a non-terminal and
$\rho(V(T'_x))<3W$. Output $V(T'_x)$ as a cluster and delete $T'_x$ from $T'$ (if $x=r$ the
tree becomes empty). When the loop stops, the remaining tree has weight less than $W$; if it
is non-empty, add its vertices to the last output cluster, whose top vertex has its parent in
the remaining tree. At least one cluster is output since $\rho q\ge W$, every cluster induces
a connected subgraph (it contains a spanning subtree), and the weights are in $[W,3W)$, or
$[W,4W)$ for the merged cluster.
\end{proof}

From now on fix $W:=2$ and assume $\rho q\ge2$. Fix the partition of
Lemma~\ref{lem:clusters}, and write $G_j:=G[V_j]$, $S_j:=X\cap V_j$, and $j(u)$ for the index
of the cluster containing $u$. Thus
\begin{equation}\label{eq:clustersize}
2/\rho\le|S_j|<8/\rho\qquad\text{for every }j.
\end{equation}
Distinct cluster graphs are vertex-disjoint, hence edge-disjoint. An edge of $G$ whose
endpoints lie in different clusters belongs to no cluster graph.

\begin{definition}
Let $s\in S_j$. A \emph{bad bridge for $s$} is an edge $e$ of $G_j$ that is a bridge of
$G_j$ but not a bridge of $G$, such that the component $Z_e(s)$ of $G_j-e$ containing $s$
satisfies $\rho|S_j\cap Z_e(s)|<1$. The terminal $s$ is \emph{good} if it has no bad bridge.
\end{definition}

\begin{lemma}\label{lem:bad}
Let $s\in S_j$.
\begin{enumerate}[label=(\roman*),nosep]
\item If $s$ is good, then $s$ spreads in $(G_j,S_j)$.
\item If $s$ is not good, there is a bad bridge $e_s$ for $s$ such that
$Z_s:=Z_{e_s}(s)\subseteq Z_e(s)$ for every bad bridge $e$ for $s$. Moreover there is an edge
$f_s=x_sy_s$ of $G$ with $x_s\in Z_s$ and $y_s\notin V_j$.
\end{enumerate}
Both can be computed in polynomial time.
\end{lemma}
\begin{proof}
(i) follows from Lemma~\ref{lem:criterion} with $Q=G_j$ and $R=S_j$, since
$\rho|S_j|\ge2$.

(ii) Let $e\ne e'$ be bad bridges for $s$, with $Z=Z_e(s)$, $Z'=Z_{e'}(s)$ and far sides
$F=V_j\setminus Z$, $F'=V_j\setminus Z'$. We claim $Z\subseteq Z'$ or $Z'\subseteq Z$. The
graph $G_j[F]$ is a component of $G_j-e$, so it is connected. If $e'$ is not an edge of
$G_j[F]$, then $F$ lies in one component of $G_j-e'$, so $F\subseteq F'$ (i.e.\ $Z'\subseteq Z$)
or $F\cap F'=\emptyset$. If $e'$ is an edge of $G_j[F]$, then $G_j[Z]$ is connected, contains
$s$ and avoids $e'$, so $Z\subseteq Z'$. It remains to exclude $F\cap F'=\emptyset$: then
$F'\subseteq Z$, and since $\rho|S_j\cap Z'|<1$ and $\rho|S_j|\ge2$ we get
$\rho|S_j\cap Z|\ge\rho|S_j\cap F'|>1$, contradicting that $e$ is bad. Hence the sets $Z_e(s)$
over bad bridges $e$ form a chain; distinct bridges have distinct sides (as
$\delta_{G_j}(Z_e(s))=\{e\}$), and we let $e_s$ be the one with the smallest side.

Since $e_s$ is not a bridge of $G$ and $V_j\setminus Z_s\neq\emptyset$, the cut
$\delta_G(Z_s)$ contains an edge $f_s\ne e_s$; write $f_s=x_sy_s$ with $x_s\in Z_s$. If
$y_s\in V_j$, then $f_s$ would be an edge of $G_j$ joining the two components of $G_j-e_s$, so
$f_s=e_s$; hence $y_s\notin V_j$.
\end{proof}

\begin{lemma}[tagging]\label{lem:tag}
Let $s\in S_j$ be not good, and let $A$ be a connected subgraph of $G$ that is edge-disjoint
from $G_j$, contains $x_s$, and contains $S_{j'}$ for some $j'\ne j$. Then $s$ spreads in
$(G_j\cup A,\ S_j\cup S_{j'})$.
\end{lemma}
\begin{proof}
Let $Q=G_j\cup A$ (connected) and $R=S_j\cup S_{j'}$, so $\rho|R|\ge2$. We verify the
hypothesis of Lemma~\ref{lem:criterion}. Let $e$ be a bridge of $Q$ that is not a bridge of
$G$, let $Z$ be the component of $Q-e$ containing $s$, and let $Z_0=Z\cap V_j$. Every edge of
$G_j$ between $Z_0$ and $V_j\setminus Z_0$ lies in $\delta_Q(Z)=\{e\}$. If there is no such
edge, connectivity of $G_j$ gives $V_j\subseteq Z$ and $\rho|R\cap Z|\ge\rho|S_j|\ge2$.
Otherwise $e\in E(G_j)$ is a bridge of $G_j$ and $Z_0=Z_e(s)$. If
$\rho|S_j\cap Z_0|\ge1$ we are done. If not, $e$ is a bad bridge for $s$, so
$x_s\in Z_s\subseteq Z_0\subseteq Z$ by Lemma~\ref{lem:bad}(ii). Since $A$ is connected,
contains $x_s$ and does not contain $e$, we get $V(A)\subseteq Z$, so
$\rho|R\cap Z|\ge\rho|S_{j'}|\ge2$.
\end{proof}

\begin{lemma}[pairing in trees; cf.~{\cite[Lemma~3.14]{CKS13}}]\label{lem:pairing}
Let $H$ be a tree and $Y$ a multiset of vertices of $H$ of even size. Then one can compute
$|Y|/2$ pairwise edge-disjoint paths in $H$ (a path may consist of a single vertex) whose
endpoints are exactly the elements of $Y$, each element used once.
\end{lemma}
\begin{proof}
Induction on $|V(H)|$. Two copies of the same vertex can be paired by a trivial path, which
uses no edge; so we may assume that every vertex occurs at most once in $Y$. If $|V(H)|=1$
then $Y=\emptyset$. Otherwise let $\ell$ be a leaf of $H$ with neighbour $u$. If
$\ell\notin Y$, apply induction to $H-\ell$. If $\ell\in Y$, let $Y'$ be the multiset obtained
from $Y$ by replacing $\ell$ by an (additional) copy of $u$, apply induction to $Y'$ in
$H-\ell$, and extend the path that ends at this copy of $u$ by the edge $u\ell$ (a trivial path
at $u$ becomes the path $u\ell$). The edge $u\ell$ is used by exactly one path.
\end{proof}

\subsection{Selecting the pairs: proof of Theorem~\ref{thm:round}}\label{sec:selection}

If $\rho q<2$, i.e.\ $\alpha q<8$, we route a single pair along a path; this suffices since
$|M|\le q/2$ gives $\alpha|M|/2^{15}<1$. Assume $\rho q\ge2$ and fix the clusters as above.

\paragraph{Step 1: one pair per cluster.}
Scan the pairs of $M$ in any order and add $st$ to a set $Z$ if neither $j(s)$ nor $j(t)$ is
marked; then mark $j(s)$ and $j(t)$.

\begin{claim}\label{cl:Z}
$|Z|\ge\rho|M|/16$, and every cluster contains terminals of at most one pair of $Z$.
\end{claim}
\begin{proof}
A pair of $M\setminus Z$ has an endpoint in a cluster marked by an earlier pair of $Z$.
Every pair of $Z$ marks two clusters with fewer than $16/\rho$ terminals in total
by~\eqref{eq:clustersize}, so fewer than $16/\rho$ pairs of $M$ (including itself) have an
endpoint there. Hence $|M|<16|Z|/\rho$. If a cluster contained terminals of two pairs of
$Z$, the later one would have been added while the cluster was marked.
\end{proof}

Let $Z_1\subseteq Z$ be the pairs with $j(s)=j(t)$. If $|Z_1|\ge|Z|/2$, route every
$st\in Z_1$ along a path in $G_{j(s)}$; these paths are edge-disjoint by
Claim~\ref{cl:Z}, and $|Z_1|\ge\rho|M|/32$. Otherwise let $Z'=Z\setminus Z_1$, so
$|Z'|\ge|Z|/2$. For a terminal $s$ of a pair of $Z'$ we call $j(s)$ its \emph{own cluster};
by Claim~\ref{cl:Z}, distinct terminals of pairs of $Z'$ have distinct own clusters, and an
own cluster contains no other terminal of a pair of $Z$.

\paragraph{Step 2: tags and an independent set.}
For every terminal $s$ of a pair of $Z'$ that is not good (in its own cluster),
Lemma~\ref{lem:bad} gives $f_s=x_sy_s$; let $\tau(s):=j(y_s)\ne j(s)$ be the cluster
\emph{tagged} by $s$. Let $D$ be the digraph on $Z'$ with an arc from $P$ to $P'\ne P$
whenever a non-good terminal $s$ of $P$ has $\tau(s)=j(u)$ for a terminal $u$ of $P'$. Every
pair has out-degree at most two, because a cluster contains terminals of at most one pair of
$Z$.

\begin{claim}\label{cl:I}
One can compute $I\subseteq Z'$ with $|I|\ge|Z'|/5$ such that no arc of $D$ joins two pairs
of $I$.
\end{claim}
\begin{proof}
Every subset $Z''\subseteq Z'$ induces at most $2|Z''|$ arcs, so it contains a pair with at most
four neighbours (in- or out-) within $Z''$. Repeatedly put such a pair into $I$ and delete it
and its neighbours.
\end{proof}

Let $I_2\subseteq I$ be the pairs $st$ with $\tau(s)=j(t)$ or $\tau(t)=j(s)$, and
$I_3=I\setminus I_2$. Suppose $|I_2|\ge|I|/2$. For $st\in I_2$ with, say, $\tau(s)=j(t)$, take
a walk consisting of a path from $s$ to $x_s$ in $G_{j(s)}$, the edge $f_s$, and a path from
$y_s$ to $t$ in $G_{j(t)}$. All edges of this walk have both endpoints in
$V_{j(s)}\cup V_{j(t)}$, and for distinct pairs of $I_2$ these unions are disjoint. So the
walks are edge-disjoint, $I_2$ is routable, and $|I_2|\ge|I|/2\ge\rho|M|/320$. From now on
assume $|I_3|\ge|I|/2$.

\begin{claim}\label{cl:transit}
If $s$ is a non-good terminal of a pair of $I_3$, then the cluster $\tau(s)$ contains no
terminal of any pair of $I$.
\end{claim}
\begin{proof}
Let $st\in I_3$. We have $\tau(s)\ne j(s)$, and $\tau(s)\ne j(t)$ since $st\notin I_2$. If
$\tau(s)=j(u)$ for a terminal $u$ of a pair $P'\in I\setminus\{st\}$, then $D$ has an arc from
$st$ to $P'$, contradicting Claim~\ref{cl:I}.
\end{proof}

The clusters $\tau(s)$ of Claim~\ref{cl:transit} are called \emph{transit clusters}. They
contain no terminal of $I$, so a transit cluster is nobody's own cluster. For a transit
cluster $C$ let $T(C)$ be the set of non-good terminals $s$ of pairs of $I_3$ with $\tau(s)=C$,
and $d(C)=|T(C)|$.

The tag edges are pairwise distinct and belong to no cluster graph: $f_s$ joins the own
cluster $j(s)$ to a transit cluster, so its endpoint in a non-transit cluster is $x_s$; if
$f_s=f_u$ for terminals $s\ne u$ of pairs of $I_3$, then $x_s=x_u$, so $j(s)=j(u)$, which is
impossible for distinct terminals.

\paragraph{Step 3: pairing through transit clusters.}
For every transit cluster $C$ with $d(C)\ge2$, fix a spanning tree $H_C$ of $G_C$; if $d(C)$ is
odd, declare one terminal of $T(C)$ \emph{isolated}. Apply Lemma~\ref{lem:pairing} to the
multiset $\{y_s:s\in T(C)\text{ not isolated}\}$ in $H_C$; if a path $P$ pairs $y_s$ with
$y_{s'}$, call $s$ and $s'$ \emph{partners} and set $P_{ss'}:=P$. Let $B$ be the graph on the
terminals whose edges join partners; every terminal has degree at most one in $B$.

\begin{claim}\label{cl:Y}
The set $Y$ of pairs of $I_3$ with no isolated terminal satisfies $|Y|\ge|I_3|/3$.
\end{claim}
\begin{proof}
A transit cluster has at most one isolated terminal, and only if $d(C)\ge3$. Since every
terminal lies in at most one set $T(C)$, the number of isolated terminals is at most
$\frac13\sum_Cd(C)\le\frac23|I_3|$, and each of them excludes one pair.
\end{proof}

\paragraph{Step 4: sparsification.}
Two distinct pairs of $Y$ \emph{conflict} if a terminal of one is a partner of a terminal of
the other. Each pair conflicts with at most two others, so greedily we find a conflict-free
$Y'\subseteq Y$ with $|Y'|\ge|Y|/3$. Let $Y'_a$ be the pairs $st\in Y'$ in which $s$ and $t$
are partners, and $Y'_b=Y'\setminus Y'_a$.

Suppose $|Y'_a|\ge|Y'|/2$. For $st\in Y'_a$ partnered through $C$, take the walk formed by a
path from $s$ to $x_s$ in $G_{j(s)}$, the edge $f_s$, the path $P_{st}$ in $H_C$, the edge
$f_t$, and a path from $x_t$ to $t$ in $G_{j(t)}$. Own clusters are distinct and are not
transit clusters, pairing paths in one tree $H_C$ are edge-disjoint, distinct transit clusters
are vertex-disjoint, and tag edges are distinct and lie in no cluster graph. Hence these walks
are edge-disjoint and $Y'_a$ is routable.

\paragraph{Step 5: spreading for $Y'_b$.}
Now suppose $|Y'_b|\ge|Y'|/2$.

\begin{claim}\label{cl:cases}
Let $s$ be a terminal of a pair $st\in Y'_b$. Then exactly one of the following holds:
($\alpha$) $s$ is good; ($\beta$) $s$ is not good and $d(\tau(s))=1$; ($\gamma$) $s$ is not
good, $d(\tau(s))\ge2$, and $s$ has a partner $s'$ that is not a terminal of any pair of $Y'$.
\end{claim}
\begin{proof}
If $s$ is not good and $d(\tau(s))\ge2$, then $s$ is not isolated (as $st\in Y$), so it has a
partner $s'$. We have $s'\ne t$ because $st\notin Y'_a$, and if $s'$ were a terminal of a pair
$P'\in Y'$, then $P'\ne st$ and $st$ would conflict with $P'$.
\end{proof}

For every terminal $s$ of a pair of $Y'_b$ define a subgraph $Q(s)$ and a receiver set $R(s)$:
\begin{itemize}[nosep]
\item[($\alpha$)] $Q(s)=G_{j(s)}$ and $R(s)=S_{j(s)}$;
\item[($\beta$)] with $C=\tau(s)$: $Q(s)=G_{j(s)}\cup\{f_s\}\cup G_C$ and
$R(s)=S_{j(s)}\cup S_C$;
\item[($\gamma$)] with $C=\tau(s)$: $Q(s)=G_{j(s)}\cup\{f_s\}\cup P_{ss'}\cup\{f_{s'}\}\cup
G_{j(s')}$ and $R(s)=S_{j(s)}\cup S_{j(s')}$. Here $j(s')$ is the own cluster of $s'$; we call it
the \emph{donor cluster} of $s$.
\end{itemize}

\begin{claim}\label{cl:spread}
Every such $s$ spreads in $(Q(s),R(s))$.
\end{claim}
\begin{proof}
Case ($\alpha$) is Lemma~\ref{lem:bad}(i). In case ($\beta$) apply Lemma~\ref{lem:tag} with
$A=\{f_s\}\cup G_C$: it is connected ($f_s$ joins $x_s$ to $y_s\in V_C$), contains $x_s$ and
$S_C$, and is edge-disjoint from $G_{j(s)}$ because $V_C\cap V_{j(s)}=\emptyset$. In case
($\gamma$) apply Lemma~\ref{lem:tag} with $A=\{f_s\}\cup P_{ss'}\cup\{f_{s'}\}\cup G_{j(s')}$:
it is connected (the walk $x_s,y_s,P_{ss'},y_{s'},x_{s'}$ ends in $V_{j(s')}$), contains $x_s$
and $S_{j(s')}$, and $j(s')\ne j(s)$ since $s'\notin\{s,t\}$. It is edge-disjoint from
$G_{j(s)}$: $f_s$ and $f_{s'}$ have an endpoint in $V_C$, $P_{ss'}\subseteq G_C$, and
$V_{j(s')}\cap V_{j(s)}=\emptyset$.
\end{proof}

\begin{claim}\label{cl:disjoint}
For distinct terminals $s\ne u$ of pairs of $Y'_b$, the graphs $Q(s)$ and $Q(u)$ are
edge-disjoint and $R(s)\cap R(u)=\emptyset$.
\end{claim}
\begin{proof}
Each $Q(s)$ consists of: its own cluster graph; at most two tag edges; in case ($\beta$) a
transit cluster graph $G_C$ with $d(C)=1$; in case ($\gamma$) a pairing path inside a transit
cluster with $d(C)\ge2$ and a donor cluster graph. We check that no part is shared. Own clusters
of distinct terminals are distinct. A donor cluster $j(s')$ contains the terminal $s'$ of a
pair of $I$ that is not in $Y'$; it is therefore not the own cluster of any terminal of $Y'_b$
(an own cluster contains terminals of only one pair of $Z$) and not a transit cluster. The
donors of distinct terminals are distinct (partners are unique), and the donor clusters of
distinct donors $s',u'$ are distinct: otherwise $s',u'$ would be the two terminals of one pair
of $Z'$, which lie in different clusters. A transit cluster $C$ with $d(C)=1$ is used, in case
($\beta$), only by the unique terminal of $T(C)$, and it carries no pairing path. Pairing
paths in one tree are edge-disjoint, and $P_{ss'}$ is used only by $s$, since its other end
$s'$ is not a terminal of $Y'$. Tag edges are distinct, lie in no cluster graph, and $f_{s'}$
is used only by the unique partner $s$ of $s'$. Finally, $R(s)$ is the union of $S_j$ over the
own cluster and the transit or donor cluster of $s$; we showed that these clusters are
distinct for distinct terminals, and the sets $S_j$ are disjoint.
\end{proof}

For every terminal $s$ of a pair of $Y'_b$ let $g_s$ be the flow given by
Claim~\ref{cl:spread}. By Claim~\ref{cl:disjoint}, the total load of these flows is at most
$1/2$ on every non-bridge and at most one on every bridge, and every $v\in X$ receives at
most $\rho$, from at most one terminal. Thus they form a spreading system, and
Lemma~\ref{lem:assembly} routes $Y'_b$ with congestion one.

\paragraph{Counting.}
In every branch the routed set has size at least
$|Y'|/2\ge|Y|/6\ge|I_3|/18\ge|I|/36\ge|Z'|/180\ge|Z|/360\ge\rho|M|/5760=\alpha|M|/23040$,
and $23040<2^{15}$; the other branches give more. If the routed set happens to be empty
(possible only when $\alpha|M|/2^{15}<1$), we route a single pair instead. Every step (clusters, bridges, tags, greedy
selections, pairing, maximum flows, the final linear program and cycle cancelling) takes
polynomial time. By Lemma~\ref{lem:reduce}(d) the routing transfers to the original graph.
This proves Theorem~\ref{thm:round}.\qed

\subsection{Non-uniform paired weights: proof of Corollary~\ref{cor:paired}}\label{sec:paired}

If $\pi(X)<2^{19}$, some pair has positive weight, so its endpoints are connected (the
product demand between them is positive); route that pair alone. Assume $\pi(X)\ge 2^{19}$
and let $q=|X|$ and $\phi_0=\pi(X)/(4q)$. Discard all pairs with $\pi(s)<\phi_0$; their total
weight is less than $q\phi_0=\pi(X)/4$. For each remaining terminal $u$ let
$n(u)=\lfloor\pi(u)/\phi_0\rfloor\ge1$ and $\pi'(u)=\phi_0n(u)\ge\pi(u)/2$, and $\pi'=0$ on
discarded terminals. Then $\pi'(X)\ge\frac38\pi(X)$ and $n(u)\le1/\phi_0\le 4q$. Let
$\lambda=\pi'(X)/\pi(X)\in[\frac38,1]$ and $\pi''=\lambda\pi'$. Since
$\pi''(u)\pi''(v)/\pi''(X)=\pi'(u)\pi'(v)/\pi(X)\le\pi(u)\pi(v)/\pi(X)$, the set
$X''=\mathrm{supp}(\pi'')$ is $\pi''$-flow-linked, and $\pi''(X)\ge\frac{9}{64}\pi(X)$.

Construct $G'$ from $G$ by attaching, to every $u\in X''$, $n(u)$ new pendant vertices
(\emph{copies} of $u$), as in~\cite[Lemma~1.1]{CKS05}. For each remaining pair $st$ (note $n(s)=n(t)$) match the $r$-th copy
of $s$ with the $r$-th copy of $t$; let $M_c$ be the resulting $q'=\sum_u n(u)$ copies'
matching. Put $\phi=\lambda\phi_0\le\frac14$, so that $\pi''(u)=\phi\,n(u)$ and
$\pi''(X)=\phi q'$. We claim the set of copies is $\phi$-flow-linked in $G'$. Between copies of
distinct $u,v$ the total demand is $n(u)n(v)\phi/q'=\pi''(u)\pi''(v)/\pi''(X)$, which we route
along the $\pi''$-product flow between $u$ and $v$ (at most the $\pi$-product flow) extended by
the pendant edges; between two copies of the same $u$ we route through $u$. Each pendant edge
carries at most $\phi(q'-1)/q'<1$, and the original edges carry at most one. All copies lie in
one component of $G'$, to which we apply Theorem~\ref{thm:round} with $\alpha=\phi$. It
returns at least $\phi|M_c|/2^{15}=\pi''(X)/2^{16}\ge 9\pi(X)/2^{22}\ge\pi(X)/2^{19}$ copy pairs
routable with congestion one.

No two of them are copies of the same pair $st$: the unique edge of $G$ at $s$ separates $s$
and all its copies from $t$ and all its copies, so two such copy pairs would each send one
unit across this edge. Deleting the pendant edges from the routing (after cycle cancelling,
a path from a copy of $s$ to a copy of $t$ is a path from $s$ to $t$ extended by two pendant
edges) gives a routing of the corresponding pairs of $M$ in $G$ with congestion one. All
quantities are polynomial since $n(u)\le4q$.\qed

\section{Well-linked decomposition with near-logarithmic loss}\label{sec:decomp}

\begin{theorem}\label{thm:decomp}
Let $G$ be a graph and $(s_i,t_i)_{i\in[k]}$ pairs whose $2k$ endpoints are distinct vertices
of degree one. Let $f$ be a feasible multicommodity flow routing $x_i\in[0,1]$ between $s_i$
and $t_i$, with $F=\sum_ix_i>0$. Then one can compute in polynomial time, using a constant
number of calls to the algorithm of Theorem~\ref{thm:bjs} on $G$ with an integral weighting of
total weight at most $4k$, disjoint sets $U_1,\dots,U_\ell\subseteq V(G)$, and for every $j$ a
set $P_j$ of pairs with both endpoints in $U_j$ and weights $\pi_j$ on the endpoint set $X_j$
of $P_j$, such that:
\begin{enumerate}[label=(\roman*),nosep]
\item $\pi_j(s_i)=\pi_j(t_i)\le1/8$ for $i\in P_j$; moreover all values of all $\pi_j$ are
positive integer multiples of a common rational $\phi>0$;
\item every vertex of $X_j$ has degree one in $G[U_j]$, and $X_j$ is $\pi_j$-flow-linked in
$G[U_j]$;
\item $\sum_j\pi_j(X_j)\ge 3F/(128\,\Gamma(4k))$.
\end{enumerate}
\end{theorem}

\begin{proof}
\emph{Discretization.} Let $\eps=F/(2k)$. For every $i$ let $d_i=\lfloor x_i/\eps\rfloor$ if
$x_i\ge\eps$ and $d_i=0$ otherwise (this is the uniformization idea
of~\cite[Lemma~1.1]{CKS05}, at granularity $\eps$). Let $K=\sum_id_i$. Then
$K\le\sum_ix_i/\eps=2k$. Moreover $\lfloor y\rfloor\ge y/2$ for $y\ge1$, and the pairs with
$x_i<\eps$ have total value less than $k\eps=F/2$, so
\begin{equation}\label{eq:K}
K\eps\ \ge\ \tfrac12\textstyle\sum_{i:x_i\ge\eps}x_i\ \ge\ F/4 .
\end{equation}
Let $A(s_i)=A(t_i)=d_i$ for every $i$ and $A=0$ elsewhere; this is well defined because the
endpoints are distinct, and $|A|=2K\le4k$.

\emph{One decomposition.} With $\Gamma=c_0\Gamma_1$ as in Section~\ref{sec:prelim}, for
$j=0,1,2,\dots$ run the algorithm of Theorem~\ref{thm:bjs} on $(G,A)$ with
$\phi_j=\eps/(2^{j+3}\Gamma_1(2K))$, obtaining $C_j$, and stop at the first $j$ with
$|C_j|\le\eps K/4$. Since $|C_j|\le\phi_j\cdot2K\cdot c_0\Gamma_1(2K)=c_0\eps K/2^{j+2}$, the
search stops at some $j\le\lceil\log c_0\rceil$, i.e.\ after $O(1)$ calls. Let $\phi$ and $C$
be the final values. Then
\[
|C|\le\eps K/4,\qquad \phi\ \ge\ \frac{\eps}{2^{\lceil\log c_0\rceil+3}\Gamma_1(2K)}\ \ge\
\frac{\eps}{16\,\Gamma(2K)},
\]
and $A|_U$ is $\phi$-flow-expanding in $G[U]$ for every component $U$ of $G-C$. (If $c_0$ is
known, a single call with $\phi=\eps/(8\Gamma(2K))$ suffices.)

\emph{Charging.} Call pair $i$ (with $d_i\ge1$) \emph{separated} if $s_i$ and $t_i$ lie in
different components of $G-C$. If $Z$ is the component containing $s_i$, then
$\delta_G(Z)\subseteq C$, so the flow of pair $i$ in $f$ has load at least $x_i\ge\eps d_i$ on
$C$. As $f$ has load at most one on each edge (cf.~\cite[Lemma~2.2]{CKS05}),
\[
\eps\textstyle\sum_{i\ \mathrm{separated}}d_i\ \le\ \sum_{i\ \mathrm{separated}}x_i\ \le\
|C|\ \le\ \eps K/4 .
\]
Hence the pairs that are not separated and have $d_i\ge1$ (the \emph{surviving} pairs) satisfy
$\sum_{i\ \mathrm{surviving}}d_i\ge3K/4$.

\emph{The pieces.} For every component $U$ of $G-C$ that contains a surviving pair, let $P_U$
be the set of surviving pairs with both endpoints in $U$, $X_U$ their endpoints, and
$\pi_U=\phi A|_{X_U}$. Property (i) holds since $\phi\le\eps/(8\Gamma_1(2K))\le\eps/8$ gives
$\phi d_i\le\phi x_i/\eps\le1/8$, and all values are multiples of $\phi$. For a surviving pair $i$ in $U$, the vertices $s_i\ne t_i$ are
connected in $G-C$, so the unique edge at $s_i$ is not in $C$ and the neighbour of $s_i$ lies in
$U$; thus $s_i$ (and likewise $t_i$) has degree one in $G[U]$. By Lemma~\ref{lem:monotone},
applied to the $\phi$-flow-expanding weighting $A|_U$ in $G[U]$ and $Y=X_U$, the set $X_U$ is
$\pi_U$-flow-linked in $G[U]$. This proves (ii). Distinct components are disjoint. Finally,
by~\eqref{eq:K},
\[
\sum_U\pi_U(X_U)=2\phi\sum_{i\ \mathrm{surviving}}d_i\ \ge\ \tfrac32\phi K\ \ge\
\frac{3K\eps}{32\,\Gamma(2K)}\ \ge\ \frac{3F}{128\,\Gamma(4k)},
\]
using monotonicity of $\Gamma$. All steps are polynomial.
\end{proof}

\begin{remark}[other decomposition algorithms]\label{rem:oracle}
The proof uses only the following property of Theorem~\ref{thm:bjs}: for the integral
weighting $A$ with $|A|\le4k$ and any $\phi>0$, the algorithm returns $C$ with
$|C|\le\phi|A|\gamma$ such that $A|_U$ is $\phi$-flow-expanding in $G[U]$ for every component
$U$ of $G-C$. Any algorithm with this property, for a known $\gamma$, yields
Theorem~\ref{thm:decomp} with $16\Gamma(4k)$ replaced by $8\gamma$ (take $\phi=\eps/(8\gamma)$).
\end{remark}

\section{The approximation algorithm}\label{sec:main}

\begin{proof}[Proof of Theorem~\ref{thm:main}]
Normalize the instance (Section~\ref{sec:prelim}) and solve the LP~\eqref{eq:lp} in its
compact edge form; let $x$ and $f$ be an optimal solution. If $\LP=0$, no pair is connected and
we return the empty set. Otherwise $\LP\ge\OPT\ge1$, since a single connected pair is routable.
Apply Theorem~\ref{thm:decomp} to $f$ with $F=\LP$. For every piece $(U_j,P_j,\pi_j)$ apply
Corollary~\ref{cor:paired} in $G[U_j]$; its hypotheses hold by Theorem~\ref{thm:decomp}(i)--(ii).
It returns a set of pairs of $P_j$ of size at least $\pi_j(X_j)/2^{19}$ routable in $G[U_j]$
with congestion one. The sets $U_j$ are disjoint, so the induced subgraphs are edge-disjoint and
the union of these routings is a routing of congestion one in $G$. Its size is at least
\[
\sum_j\frac{\pi_j(X_j)}{2^{19}}\ge\frac{3\,\LP}{2^{26}\,\Gamma(4k)} .
\]
(If no piece exists, we return one connected pair; Corollary~\ref{cor:paired} returns at least
one pair per piece, so the output is never empty.) All steps other than the calls to the
algorithm of~\cite{BJS26} are deterministic and polynomial. By~\eqref{eq:gamma},
$\Gamma(4k)\le C'\log(k+2)\exp(C'\sqrt{\log\log(k+4)})$ for an absolute constant $C'$, and
$\LP\ge\OPT$.
\end{proof}

\begin{corollary}[integrality gap]
For every unit-capacity instance with $k$ pairs, $\LP\le O(\Gamma(4k))\cdot\OPT
=\log^{1+o(1)}k\cdot\OPT$.
\end{corollary}

\begin{remark}[integer capacities]
If edges have integer capacities $u(e)$ and demands are one, one can replace every edge by
$u(e)$ parallel unit edges. To keep the instance polynomial, first use the argument
of~\cite[Section~2]{CKS13}: decompose an optimal edge flow into polynomially many paths and
take a basic optimal solution of the path formulation restricted to these paths. A pair with
$x_i\in(0,1)$ has a non-tight pair constraint, so the number of such pairs is at most the
number of tight edge constraints, i.e.\ at most $m$. If the pairs with $x_i=1$ carry at least
half of $\LP$, they form a routable set; otherwise restrict to the at most $m$ fractional pairs,
whose total load on any edge is at most $m$, and cap the capacities at $m$. This loses a factor
two, and Theorem~\ref{thm:main} then applies with the same asymptotic ratio.
\end{remark}

\begin{remark}[weighted version, with an extra logarithm]\label{rem:weighted}
For pairs with profits $w_i\ge0$ and the weighted relaxation $\LP_w=\max\sum_iw_ix_i$, the
standard class argument gives ratio $O(\Gamma(4k)\log k)=\log^{2+o(1)}k$. Let $W$ be the
largest profit of a pair with connected endpoints, so $\LP_w\ge W$ (pairs with disconnected
endpoints have $x_i=0$). Pairs with $w_i<W/k^2$ contribute at most $W/k\le\LP_w/k$ and are
discarded; the others fall into $O(\log k)$ classes $\{i:w_i\in(W2^{-j-1},W2^{-j}]\}$, one of
which carries at least an $\Omega(1/\log k)$ fraction of $\LP_w$. Theorem~\ref{thm:decomp} and
Corollary~\ref{cor:paired} apply to the optimal flow restricted to that class (any feasible flow
may be used), and every routed pair has profit at least half the maximum profit in its class. We
do not know how to avoid the extra logarithm; see Section~\ref{sec:discussion}.
\end{remark}

\begin{remark}[Anand's decomposition; conditional]\label{rem:anand}
Anand~\cite{Anand26} improves flow-expander decomposition to overhead $O(\log n\log\log n)$.
His node-weighted theorem (Theorem~5.1 of~\cite{Anand26}) gives overhead
$O(\log|A|\log\log|A|)$ for rational weightings with $A(v)\ge1$ for \emph{every} vertex $v$, so
that $|A|\ge n$. Our weighting vanishes outside the $2k$ terminals, so this bound does not give
an overhead in terms of $k$. For weightings that may vanish, Section~5 and Appendix~C
of~\cite{Anand26} state overhead $O(\log n\log\log n)$, with $n$ the number of vertices; we
have not verified that appendix. \emph{If} that statement holds, Remark~\ref{rem:oracle} and
Corollary~\ref{cor:paired} give a polynomial-time $O(\log n\log\log n)$-approximation for
AN-MCF (the normalized graph has $n+2k$ vertices, and by the previous remark we may assume
$k\le m$). This bound is incomparable with Theorem~\ref{thm:main} and better when
$n\le k^{O(1)}$. Obtaining $O(\log k\log\log k)$ would require a node-weighted decomposition
whose overhead depends only on the total weight even when most vertices have weight zero, or a
reduction eliminating zero-weight vertices without loss; we leave this open.
\end{remark}

\section{Lower bounds and open problems}\label{sec:discussion}

\paragraph{Lower bounds.}
Andrews, Chuzhoy, Guruswami, Khanna, Talwar and Zhang~\cite{ACGKTZ10} show that AN-MCF is hard
to approximate within $\log^{1/2-\delta}n$ unless
$\mathrm{NP}\subseteq\mathrm{ZPTIME}(n^{\mathrm{polylog}\, n})$, and that the integrality gap
of~\eqref{eq:lp} is $\log^{1/2-\delta}n$ (as quoted in~\cite{CKS13}). These bounds are stated in
terms of the number $n$ of vertices. Pairs with $s_i=t_i$ count equally in $\LP$ and $\OPT$ and
can be removed; afterwards, keeping at most $\deg(u)$ copies of each repeated pair $\{u,v\}$
changes neither $\LP$ nor $\OPT$, so one may assume $k\le n^2m$. For instances of polynomial
size ($m\le n^{O(1)}$, as produced by such reductions) this gives $\log k=O(\log n)$, and the
lower bounds read $\Omega(\log^{1/2-\delta}k)$. Thus on such instances the integrality gap of
the natural relaxation lies between $\log^{1/2-\delta}k$ and $\log^{1+o(1)}k$. The
$\Omega(\log n)$ lower bound on the overhead of expander decompositions~\cite{BJS26} concerns the
decomposition, not AN-MCF, and we do not know whether $\Omega(\log k)$ is a barrier for
approaches based on flow-linked decompositions.

\paragraph{Weighted version.}
For pairs with profits $w_i$, the decomposition step charges separated pairs by the amount of
flow crossing the cut, which bounds their total \emph{flow}; a cut of the allowed size may
separate all of the most profitable pairs. Weighting the terminals by profit makes the cut
bound grow with the profit range, and grouping pairs into profit classes loses a further
logarithm, which gives the $\log^{2+o(1)}k$ bound of Remark~\ref{rem:weighted}. We do not know
how to obtain a $\log^{1+o(1)}k$ ratio for weighted AN-MCF. (We
expect the rounding of Section~\ref{sec:rounding} to extend to weights along the lines
indicated in~\cite{CKS13}; the obstacle is the decomposition.)

\paragraph{Node capacities.}
For node-capacitated AN-MCF, \cite{CKS05} obtain polylogarithmic ratios with congestion
$1+\eps$. An analogue of Theorem~\ref{thm:main} would require vertex flow-expander
decompositions with near-logarithmic overhead for integral weightings, which~\cite{BJS26}
list as open, and a node analogue of the rounding.

\paragraph{Other questions.}
Can the ratio be improved to $O(\log k)$, or even to $O(\log^{1/2+o(1)}k)$? Can the
decomposition of Theorem~\ref{thm:decomp} improve other applications of well-linked
decompositions in general graphs, such as those in~\cite{CL16}?

\paragraph{Computational sanity checks.}
We implemented the rounding of Section~\ref{sec:rounding} (degree reduction, clusters, bad
bridges, tags, Steps~1--5, spreading flows by maximum flow) and ran it on $270$ random instances
with up to $68$ terminals, checking every intermediate assertion of the proof (nesting of bad bridges,
existence of tag edges, transit clusters free of terminals of $I$, edge-disjointness of the
graphs $Q(s)$ and of the direct routings, the spreading-system conditions) and verifying by linear
programming that each selected set is routable with congestion one. Lemmas~\ref{lem:criterion},
\ref{lem:bad} and~\ref{lem:tag} were also tested directly on $20000$ random small graphs. Because
of the constants, small instances mostly exercise the simpler branches of
Section~\ref{sec:selection}: the branches $Z_1$, $I_2$ and $Y'_b$ (with cases ($\alpha$) and
($\beta$)) occurred, while case ($\gamma$) and the branch $Y'_a$ did not occur (Lemma~\ref{lem:tag},
which case ($\gamma$) uses, was tested directly). These checks are diagnostics only and are not part of any
proof.

\section*{Acknowledgment}
This paper was prepared with substantial assistance from AI language models (Anthropic's
Claude), which were used for research exploration, proof drafting, and computational
verification; the author has checked the arguments and takes full responsibility for the
content.

\bibliographystyle{plain}
\bibliography{references}

\end{document}
